% !TEX output_directory = ../publica
\documentclass[twoside, 10pt, a4paper]{article}
\usepackage[utf8]{inputenc}
\usepackage[T1]{fontenc}
\usepackage[portuguese]{babel}

% --- FONTE PADRÃO (Estilo Didático Encorpado) ---
\usepackage{helvet} 
\renewcommand{\familydefault}{\sfdefault}

% --- GEOMETRIA E ESPAÇAMENTO ---
\usepackage[top=1.6cm, bottom=1.5cm, left=0.9cm, right=0.9cm, headheight=22pt, headsep=0.4cm]{geometry}
\usepackage{microtype} 
\usepackage{setspace}
\setstretch{1.0}
\usepackage{parskip}
\setlength{\parskip}{2pt}
\setlength{\parindent}{0pt}

\newlength{\espacoPosTitulo}\setlength{\espacoPosTitulo}{0.3cm}
\newlength{\espacoPreQuestao}\setlength{\espacoPreQuestao}{0.1cm}
\newlength{\espacoQuestao}\setlength{\espacoQuestao}{0.2cm}
\newlength{\espacoAlt}\setlength{\espacoAlt}{0.2cm}

\usepackage{enumitem}
\setlist{itemsep=0.4cm, topsep=0.1cm, parsep=0pt, partopsep=0pt}
\newcommand{\larguraTikz}{0.85\linewidth} 

% --- MATEMÁTICA E CORES ---
\usepackage{amsmath, amsfonts, amssymb, nccmath, mathastext}
\usepackage{xcolor}
\definecolor{indigo}{HTML}{4F46E5}
\definecolor{CorTitUm}{HTML}{FF6F61}
\definecolor{CorTitDois}{HTML}{FF6F61}
\definecolor{CorTitTres}{HTML}{658894}
\definecolor{CorLinha}{HTML}{B0B0B0} 
\definecolor{metal}{HTML}{7F8C8D}
\definecolor{vidro}{HTML}{3498DB}
\definecolor{ouro}{HTML}{F1C40F}
\definecolor{concreto}{HTML}{95A5A6}
\definecolor{madeira}{HTML}{D35400}
\definecolor{CinzaEscuro}{HTML}{4A4A4A} 
\definecolor{CinzaMedio}{HTML}{848484} 
\definecolor{Cinzablack}{HTML}{353535} 

% --- MINI-CABEÇALHO E RODAPÉ AUTORAL (TWOSIDE) ---
\usepackage{fancyhdr}
\pagestyle{fancy}
\fancyhf{} 
\renewcommand{\headrulewidth}{0.3pt} 
\renewcommand{\footrulewidth}{0.3pt}
\fancyhead[LE]{\small\sffamily\bfseries\color{gray!75} SISTEMA SOL-TERRA-LUA}
\fancyhead[RO]{\small\sffamily\color{gray!75} \texttt{www.profraphaelpaulino.com.br}}
\fancyfoot[LE]{\small \textbf{\thepage}}
\fancyfoot[RO]{\small \textbf{\thepage}}

% --- CAIXAS DE DESTAQUE ---
\usepackage[most]{tcolorbox}
\newtcolorbox{boxresponda}{colback=white, colframe=gray!45, boxrule=0pt, leftrule=1.7pt, sharp corners, enlarge left by=0.4cm, width=\linewidth-0.4cm, left=8pt, right=0pt, top=2pt, bottom=2pt, fontupper=\small}
\definecolor{CorDicaBorda}{HTML}{17c1be} 
\definecolor{CorDicaFundo}{HTML}{f3fbfb} 
\newtcolorbox{boxdica}{colback=CorDicaFundo, colframe=CorDicaBorda, boxrule=0pt, leftrule=2.5pt, sharp corners, width=\linewidth, left=8pt, right=8pt, top=6pt, bottom=6pt, halign=center, fontupper=\small}

% --- TÍTULOS ---
\usepackage{titlesec}
\titleformat{\section}{\color{black}\large\bfseries}{\thesection}{0em}{\vspace{0.1cm}\titlerule[0.8pt]}
\titlespacing*{\section}{0pt}{10pt}{6pt} 
\titleformat{\subsubsection}[runin]{\color{black}\bfseries}{}{0em}{}
\titlespacing*{\subsubsection}{0pt}{0pt}{0.15cm}

% --- COLUNAS E GRÁFICOS ---
\usepackage{multicol}
\setlength{\columnsep}{1cm}
\setlength{\columnseprule}{0.5pt}
\def\columnseprulecolor{\color{CorLinha}}
\usepackage{adjustbox, tikz, pgfplots, circuitikz}
\usetikzlibrary{shadows, shapes.misc, positioning, calc}
\pgfplotsset{compat=1.18}

\begin{document}
\thispagestyle{empty}
\color{Cinzablack}
\vspace*{-1.8cm}

% --- CABEÇALHO PRINCIPAL AUTORAL (Apenas 1ª página) ---
\noindent
\begin{minipage}{\textwidth}
    \fontfamily{phv}\selectfont
    \noindent
    \begin{minipage}[t]{0.70\linewidth}
        {\Large \bfseries \MakeUppercase{Caderno de Atividades Suplementares}}\\[0.1cm]
        {\large \textmd{\textcolor{CinzaEscuro}{Unidade: SISTEMA SOL-TERRA-LUA}}}
    \end{minipage}%
    \begin{minipage}[t]{0.30\linewidth}
        \raggedleft
        \small \textbf{Prof. Raphael Paulino}\\
        \textsf{\small \textcolor{indigo}{\texttt{profraphaelpaulino.com.br}}}
    \end{minipage}
    
    \vspace{0.15cm}
    \noindent\rule{\linewidth}{1.2pt}
    \vspace{-0.1cm}
    \footnotesize\textsf{\textbf{ÁREA:} FÍSICA \hfill \textbf{NÍVEL:} EF II (8º Ano)}
    \vspace{0.1cm}
    \noindent\rule{\linewidth}{0.4pt}
\end{minipage}

\par\vspace{0.3cm} 
\raggedcolumns
\begin{multicols*}{2}
\vspace{0.1cm}
\subsection*{Capítulo 1: O Sistema Sol-Terra-Lua}
\vspace{-0.2cm}\noindent\rule{\linewidth}{1.5pt} 
\par\vspace{\espacoPosTitulo}
\subsubsection*{1.} Uma feira de ciências montou um modelo em escala do sistema Sol-Terra-Lua. No pátio da escola, o grupo de alunos posicionou os astros respeitando a proporção de distâncias. Considere o esquema montado pelos alunos, onde a malha quadriculada no chão ajuda a identificar as posições. 

\begin{center}
% ==========================================
% CONTROLE INDIVIDUAL DESTA IMAGEM:
% 1. CORTAR BORDAS: Mude os zeros do trim (Esquerda, Baixo, Direita, Topo). Ex: trim=1cm 0cm 1cm 0cm
% 2. TAMANHO: Para encolher só esta imagem, troque \larguraTikz por um valor (Ex: 0.5\linewidth)
% ==========================================
\begin{adjustbox}{trim=0cm 0cm 0cm 0cm, clip, max width=\larguraTikz}
\begin{tikzpicture}
% --- PALETA HEXADECIMAL MODERNA ---
\definecolor{bgchao}{HTML}{E8F6F3}
\definecolor{solcolor}{HTML}{F39C12}
\definecolor{terracolor}{HTML}{2980B9}
\definecolor{luacolor}{HTML}{95A5A6}
\definecolor{linhamalha}{HTML}{BDC3C7}

% --- LAYER 1: Fundo e Malha ---
\fill[bgchao, rounded corners=4pt] (-1, -2) rectangle (9, 3);
\draw[step=1cm, linhamalha, very thin] (-1,-2) grid (9,3);

% --- APLICANDO DROP SHADOW E CAMADAS ---
\node[circle, fill=solcolor, minimum size=1.5cm, drop shadow={opacity=0.3, shadow xshift=2pt, shadow yshift=-2pt}] (Sol) at (0,0) {};
\node[fill=white, inner sep=1pt, text=solcolor, font=\bfseries\footnotesize] at (0,-1.2) {Sol};

\node[circle, fill=terracolor, minimum size=0.6cm, drop shadow={opacity=0.3, shadow xshift=2pt, shadow yshift=-2pt}] (Terra) at (6,0) {};
\node[fill=white, inner sep=1pt, text=terracolor, font=\bfseries\footnotesize] at (6,-0.8) {Terra};

\node[circle, fill=luacolor, minimum size=0.25cm, drop shadow={opacity=0.3, shadow xshift=1pt, shadow yshift=-1pt}] (Lua) at (6.5,1) {};
\node[fill=white, inner sep=1pt, text=luacolor, font=\bfseries\tiny] at (6.5, 1.4) {Lua};

% Setas de proporção (Mascaradas para não cruzar malha sujando a leitura)
\draw[<->, thick, dashed] (Sol) -- node[fill=bgchao, inner sep=2pt, font=\bfseries\footnotesize] {\textbf{d}} (Terra);
\draw[<->, thick, dashed] (Terra) -- (Lua);
\node[fill=bgchao, inner sep=1pt, font=\bfseries\tiny, align=center] at (6.8, 0.5) {\textbf{x}};
\end{tikzpicture}
\end{adjustbox}
\end{center}

Na montagem, cada quadrado da malha equivale a $ \mfrac{1}{4} $ de Unidade Astronômica (UA) no modelo. Determine a distância \textbf{d} representada no pátio da escola, em Unidades Astronômicas, e justifique por que, na vida real, a representação de \textbf{x} na mesma escala de \textbf{d} seria inviável para esse modelo visual visualizado a olho nu.
\par\vspace{\espacoQuestao}

\noindent\textcolor{CorLinha}{\rule{\linewidth}{0.5pt}} 
\par\vspace{\espacoPreQuestao}
\subsubsection*{2.} A luz viaja no vácuo a uma velocidade aproximada de \textbf{300.000 km/s}. Sabe-se que a distância média entre a Terra e o Sol é de cerca de \textbf{150.000.000 km}, enquanto a distância média da Terra à Lua é de aproximadamente \textbf{384.000 km}. Demonstre matematicamente, através do cálculo de tempo de viagem da luz (usando regras proporcionais), por que observamos a Lua praticamente em "tempo real" e o Sol com um atraso significativo.
\par\vspace{\espacoQuestao}

\noindent\textcolor{CorLinha}{\rule{\linewidth}{0.5pt}} 
\par\vspace{\espacoPreQuestao}
\subsubsection*{3.} Sobre os componentes do Sistema Sol-Terra-Lua e suas proporções, julgue as afirmativas a seguir com V (Verdadeiro) ou F (Falso) e justifique aquelas que considerar falsas:
\begin{enumerate}[label=\Roman*., itemsep=\espacoAlt]
    \item O Sol é a fonte primária de energia do sistema, sendo uma estrela que emite luz própria devido a reações nucleares, enquanto Terra e Lua são astros iluminados.
    \item O diâmetro da Terra é aproximadamente quatro vezes maior que o da Lua, o que justifica a diferença de gravidade, mas no céu Sol e Lua parecem ter tamanhos similares apenas devido à perspectiva de distância.
    \item Se a distância da Terra até a Lua dobrasse de tamanho, o tempo que a luz do Sol demora para chegar à Terra também dobraria.
\end{enumerate}
\par\vspace{\espacoQuestao}

\noindent\textcolor{CorLinha}{\rule{\linewidth}{0.5pt}} 
\par\vspace{\espacoPreQuestao}
\subsubsection*{4.} (Estilo VUNESP) O projeto astronômico "Via Láctea na Praça" propõe a criação de esferas de diferentes materiais para representar planetas. Se a Terra for representada por uma bola de gude com \textbf{1,5 cm} de diâmetro, o Sol deverá ser representado por uma esfera maior, montada em uma praça adjacente. Sabendo que o diâmetro real do Sol é cerca de \textbf{109 vezes} maior que o da Terra, calcule qual deve ser o diâmetro, em metros, da esfera que representará o Sol neste projeto escolar.
\par\vspace{\espacoQuestao}

\noindent\textcolor{CorLinha}{\rule{\linewidth}{0.5pt}} 
\par\vspace{\espacoPreQuestao}
\subsubsection*{5.} Durante uma aula de ciências, um aluno fictício chamado Marcos fez a seguinte afirmação sobre a dinâmica do nosso sistema planetário: 

\begin{boxresponda}
\textbf{Responda:} "A Terra orbita o Sol porque o Sol é o corpo mais quente do sistema, possuindo forte energia térmica que atrai a Terra. Já a Lua orbita a Terra porque ela está presa dentro da atmosfera terrestre, não conseguindo escapar para o espaço sideral."\\
\\
Diagnostique os \textbf{dois} erros conceituais (físicos e orbitais) na argumentação de Marcos e reescreva a explicação utilizando o conceito correto que define quem orbita quem no Sistema Sol-Terra-Lua.
\end{boxresponda}
\par\vspace{\espacoQuestao}

\subsection*{Capítulo 2: Movimentos da Terra}
\vspace{-0.2cm}\noindent\rule{\linewidth}{1.5pt} 
\par\vspace{\espacoPosTitulo}
\subsubsection*{6.} A inclinação do eixo de rotação terrestre é o principal fator responsável pelas características de insolação em diferentes épocas do ano. O esquema gráfico ilustra a incidência dos raios solares sobre o globo terrestre em um determinado momento da translação. 

\begin{center}
% ==========================================
% CONTROLE INDIVIDUAL DESTA IMAGEM:
% 1. CORTAR BORDAS: Mude os zeros do trim (Esquerda, Baixo, Direita, Topo). Ex: trim=1cm 0cm 1cm 0cm
% 2. TAMANHO: Para encolher só esta imagem, troque \larguraTikz por um valor (Ex: 0.5\linewidth)
% ==========================================
\begin{adjustbox}{trim=0cm 0cm 0cm 0cm, clip, max width=\larguraTikz}
\begin{tikzpicture}
% --- PALETA HEXADECIMAL MODERNA ---
\definecolor{terra}{HTML}{3498DB}
\definecolor{terraescura}{HTML}{2C3E50}
\definecolor{raios}{HTML}{F1C40F}
\definecolor{eixo}{HTML}{E74C3C}

% --- LAYER 1: Raios Solares ---
\foreach \y in {-2, -1, 0, 1, 2} {
    \draw[->, very thick, raios, drop shadow={opacity=0.1}] (-3, \y) -- (-1.2, \y);
}
\node[fill=white, text=raios, font=\bfseries] at (-2, 2.5) {Raios Solares};

% --- LAYER 2: Planeta Terra (com sombreamento para noite) ---
\begin{scope}[rotate=-23.5]
    % Lado noite (direito)
    \clip (0,-2.5) rectangle (2.5,2.5);
    \fill[terraescura, drop shadow={opacity=0.2}] (0,0) circle (2cm);
\end{scope}

\begin{scope}[rotate=-23.5]
    % Lado dia (esquerdo)
    \clip (-2.5,-2.5) rectangle (0,2.5);
    \fill[terra, drop shadow={opacity=0.2}] (0,0) circle (2cm);
    % Linha do Equador e Trópicos
    \draw[dashed, white, thick] (-2,0) -- (2,0) node[right, font=\tiny\bfseries, fill=white, text=black, inner sep=1pt] {Equador};
    \draw[dashed, white, thin] (-1.73,1) -- (1.73,1);
    \draw[dashed, white, thin] (-1.73,-1) -- (1.73,-1);
\end{scope}

% Eixo de rotação inclinado
\draw[thick, eixo] (0, -2.8) -- (0, 2.8) node[above, font=\footnotesize\bfseries, fill=white, inner sep=1pt] {Eixo de Rotação};
\draw[thick, black, dashed] (-1, -2.4) -- (1, 2.4); % Eixo vertical de referência

% Angulo
\draw[eixo, ->, thick] (0, 1.8) arc (90:113.5:1.8);
\node[font=\footnotesize\bfseries, text=eixo, fill=white, inner sep=1pt] at (-0.6, 2) {23,5$^\circ$};

\end{tikzpicture}
\end{adjustbox}
\end{center}

Analisando exclusivamente a geometria de incidência luminosa no esquema, determine em qual hemisfério (Norte ou Sul) está ocorrendo o verão e justifique sua resposta demonstrando como a área de iluminação (lado claro) se distribui em relação ao Equador.
\par\vspace{\espacoQuestao}

\noindent\textcolor{CorLinha}{\rule{\linewidth}{0.5pt}} 
\par\vspace{\espacoPreQuestao}
\subsubsection*{7.} O movimento de rotação da Terra ocorre de Oeste para Leste, com duração de aproximadamente \textbf{24 horas}. Como consequência direta desse sentido de giro, observamos o movimento aparente do Sol no céu durante o dia. Considerando duas cidades brasileiras, Recife (PE) no litoral e Rio Branco (AC) no extremo oeste do país, explique qual delas vê o nascer do Sol primeiro e como isso se relaciona com o sentido da rotação planetária.
\par\vspace{\espacoQuestao}

\noindent\textcolor{CorLinha}{\rule{\linewidth}{0.5pt}} 
\par\vspace{\espacoPreQuestao}
\subsubsection*{8.} Analise os itens abaixo relacionados ao movimento de translação e à ocorrência das estações do ano, assinalando V para verdadeiro e F para falso. Corrija matematicamente ou conceitualmente as afirmativas incorretas:
\begin{enumerate}[label=\Roman*., itemsep=\espacoAlt]
    \item O periélio (ponto da órbita mais próximo do Sol) é a causa exclusiva do verão, pois a Terra recebe mais calor do que no afélio, anulando os efeitos do eixo.
    \item Nos equinócios, os raios solares incidem perpendicularmente sobre a Linha do Equador, resultando em dias e noites com durações aproximadamente iguais em todo o planeta.
    \item Se o eixo da Terra não possuísse inclinação (fosse totalmente reto), a translação não geraria estações do ano distintas, mantendo o clima uniforme nas latitudes.
\end{enumerate}
\par\vspace{\espacoQuestao}

\noindent\textcolor{CorLinha}{\rule{\linewidth}{0.5pt}} 
\par\vspace{\espacoPreQuestao}
\subsubsection*{9.} (ENEM - Adaptada) O gráfico documenta a variação do número de horas de luz solar ao longo dos meses do ano para uma cidade X.

\begin{center}
% ==========================================
% CONTROLE INDIVIDUAL DESTA IMAGEM:
% ==========================================
\begin{adjustbox}{trim=0cm 0cm 0cm 0cm, clip, max width=\larguraTikz}
\begin{tikzpicture}
\begin{axis}[
    width=8cm, height=5cm,
    axis lines=middle,
    xmin=1, xmax=12,
    ymin=8, ymax=16,
    xlabel={\textbf{Mês}},
    ylabel={\textbf{Horas de Luz}},
    xtick={1,3,6,9,12},
    xticklabels={Jan, Mar, Jun, Set, Dez},
    ytick={10, 12, 14},
    grid=major,
    grid style={dashed, gray!30},
    thick,
    every axis plot/.append style={ultra thick}
]
% --- APLICANDO CURVA SUAVE ---
\addplot[color=CorTitUm, smooth, tension=0.7] coordinates {
    (1, 14.5)
    (3, 12)
    (6, 9.5)
    (9, 12)
    (12, 14.5)
};
\node[fill=white, inner sep=2pt, font=\bfseries] at (axis cs:6, 10.5) {Cidade X};
\end{axis}
\end{tikzpicture}
\end{adjustbox}
\end{center}

Com base na interpretação gráfica das horas de insolação e sua variação cíclica, conclua em qual hemisfério a Cidade X está localizada. Em sua justificativa, aponte explicitamente os meses que caracterizam o solstício de inverno e o solstício de verão para esta localidade.
\par\vspace{\espacoQuestao}

\noindent\textcolor{CorLinha}{\rule{\linewidth}{0.5pt}} 
\par\vspace{\espacoPreQuestao}
\subsubsection*{10.} Durante um experimento escolar de observação indireta, um aluno posiciona uma haste reta de \textbf{1 metro} verticalmente no solo plano. Ele decide medir o comprimento da sombra projetada pela haste exatamente ao meio-dia solar, durante os dias de solstício de verão e de inverno em sua cidade (localizada na região Sul do Brasil). Descreva analiticamente como será o comportamento geométrico dessa sombra: em qual dos dois eventos a sombra projetada será mais longa e por que isso ocorre?
\par\vspace{\espacoQuestao}

\noindent\textcolor{CorLinha}{\rule{\linewidth}{0.5pt}} 
\par\vspace{\espacoPreQuestao}
\subsubsection*{11.} Um jornal publicou uma matéria sobre a organização do calendário civil contendo o seguinte trecho:

\begin{boxresponda}
\textbf{Responda:} "O ano bissexto ocorre a cada quatro anos com a adição do dia 29 de fevereiro. Isso é necessário porque a Terra demora exatos 365 dias e 12 horas para dar a volta no Sol. Sem essa correção, as estações do ano se inverteriam de 6 em 6 meses."\\
\\
Utilizando o valor astronômico real da translação terrestre (aprox. 365 dias e 6 horas), refute a afirmação matemática do jornal e explique qual seria o verdadeiro impacto nas estações do ano a longo prazo se não utilizássemos o ano bissexto.
\end{boxresponda}
\par\vspace{\espacoQuestao}

\noindent\textcolor{CorLinha}{\rule{\linewidth}{0.5pt}} 
\par\vspace{\espacoPreQuestao}
\subsubsection*{12.} O monumento arqueológico de Stonehenge, na Inglaterra, possui pedras alinhadas especificamente para marcar eventos astronômicos anuais ditados pela translação e inclinação terrestre. Observe a modelagem do alinhamento da luz do Sol nascendo em relação a uma pedra de referência.

\begin{center}
% ==========================================
% CONTROLE INDIVIDUAL DESTA IMAGEM:
% ==========================================
\begin{adjustbox}{trim=0cm 0cm 0cm 0cm, clip, max width=\larguraTikz}
\begin{tikzpicture}
\definecolor{pedra}{HTML}{7F8C8D}
\definecolor{grama}{HTML}{27AE60}
\definecolor{luz}{HTML}{F1C40F}

% --- LAYER 1: Cenário ---
\fill[grama, rounded corners=2pt] (-2, -1) rectangle (6, 0);

% --- LAYER 2: Pedras (Drop Shadow) ---
\fill[pedra, drop shadow={opacity=0.3, shadow xshift=2pt, shadow yshift=-2pt}] (0,0) rectangle (1.5, 3);
\fill[pedra, drop shadow={opacity=0.3, shadow xshift=2pt, shadow yshift=-2pt}] (-1.5,0) rectangle (-0.5, 2.5);
\fill[pedra, drop shadow={opacity=0.3, shadow xshift=2pt, shadow yshift=-2pt}] (2.5,0) rectangle (3.5, 2.5);

% Lintel superior
\fill[pedra, drop shadow={opacity=0.3, shadow xshift=2pt, shadow yshift=-2pt}] (-0.2,3) rectangle (1.7, 3.5);

% --- LAYER 3: Luz ---
\draw[->, ultra thick, luz, dashed] (5, 2) -- node[above, sloped, text=black, font=\bfseries\footnotesize, fill=white, inner sep=1pt] {Nascer do Sol} (1.5, 0);
\draw[->, ultra thick, luz, dashed] (5, 2.5) -- (1.5, 0.5);

\node[fill=grama, text=white, font=\bfseries] at (4, -0.5) {Plano de Observação};
\end{tikzpicture}
\end{adjustbox}
\end{center}

Se o feixe de luz atinge perfeitamente o centro geométrico do monumento apenas no dia em que o hemisfério Norte atinge seu grau máximo de insolação anual, a qual fenômeno astronômico exato (incluindo o nome da estação) a pedra principal está alinhada?
\par\vspace{\espacoQuestao}
\subsection*{Capítulo 3: Movimentos da Lua}
\vspace{-0.2cm}\noindent\rule{\linewidth}{1.5pt} 
\par\vspace{\espacoPosTitulo}
\subsubsection*{13.} O ciclo lunar é governado pela posição relativa entre a Terra, a Lua e o Sol. O esquema representa uma vista superior do plano orbital da Lua (representada em quatro posições numéricas) ao redor da Terra.

\begin{center}
% ==========================================
% CONTROLE INDIVIDUAL DESTA IMAGEM:
% 1. CORTAR BORDAS: Mude os zeros do trim (Esquerda, Baixo, Direita, Topo). Ex: trim=1cm 0cm 1cm 0cm
% 2. TAMANHO: Para encolher só esta imagem, troque \larguraTikz por um valor (Ex: 0.5\linewidth)
% ==========================================
\begin{adjustbox}{trim=0cm 0cm 0cm 0cm, clip, max width=\larguraTikz}
\begin{tikzpicture}
\definecolor{fundo}{HTML}{2C3E50}
\definecolor{terra}{HTML}{3498DB}
\definecolor{lua}{HTML}{BDC3C7}
\definecolor{luzsol}{HTML}{F1C40F}

% LAYER 1: Fundo
\fill[fundo, rounded corners=4pt, drop shadow={opacity=0.4, shadow xshift=3pt, shadow yshift=-3pt}] (-4,-3.5) rectangle (5,3.5);

% LAYER 2: Luz do Sol
\foreach \x in {-2.5, -1, 0.5, 2} {
    \draw[->, ultra thick, luzsol] (5, \x) -- (3.5, \x);
}
\node[fill=fundo, text=luzsol, font=\bfseries] at (4.2, 3) {Sol};

% LAYER 3: Órbita e Terra
\draw[dashed, white, thick] (0,0) circle (2.2cm);
\fill[terra, drop shadow={opacity=0.5}] (0,0) circle (0.6cm);
\node[fill=white, text=terra, font=\bfseries\tiny, inner sep=1pt] at (0, -0.9) {Terra};

% LAYER 4: Luas nas Posições (Lado direito iluminado)
% Pos 1 (Direita)
\fill[black] (2.2,0) arc (-90:90:0.3cm) -- cycle;
\fill[lua] (2.2,0) arc (270:90:0.3cm) -- cycle;
\node[fill=white, text=black, font=\bfseries, inner sep=2pt, circle] at (2.8, 0.5) {1};

% Pos 2 (Cima)
\fill[black] (0,2.2) arc (-90:90:0.3cm) -- cycle;
\fill[lua] (0,2.2) arc (270:90:0.3cm) -- cycle;
\node[fill=white, text=black, font=\bfseries, inner sep=2pt, circle] at (-0.6, 2.7) {2};

% Pos 3 (Esquerda)
\fill[black] (-2.2,0) arc (-90:90:0.3cm) -- cycle;
\fill[lua] (-2.2,0) arc (270:90:0.3cm) -- cycle;
\node[fill=white, text=black, font=\bfseries, inner sep=2pt, circle] at (-2.8, -0.5) {3};

% Pos 4 (Baixo)
\fill[black] (0,-2.2) arc (-90:90:0.3cm) -- cycle;
\fill[lua] (0,-2.2) arc (270:90:0.3cm) -- cycle;
\node[fill=white, text=black, font=\bfseries, inner sep=2pt, circle] at (0.6, -2.7) {4};

\end{tikzpicture}
\end{adjustbox}
\end{center}

Baseando-se na perspectiva de um observador localizado na Terra, relacione cada posição numérica (1, 2, 3 e 4) à sua respectiva fase lunar. Em seguida, justifique por que a face iluminada da Lua na Posição 1 não pode ser vista durante a noite terrestre.
\par\vspace{\espacoQuestao}

\noindent\textcolor{CorLinha}{\rule{\linewidth}{0.5pt}} 
\par\vspace{\espacoPreQuestao}
\subsubsection*{14.} Apesar de orbitar a Terra e realizar seu próprio movimento de rotação, um observador terrestre vê sempre a mesma face da Lua, independentemente da fase em que ela se encontra. Sobre esse fenômeno dinâmico, julgue as afirmativas a seguir com V (Verdadeiro) ou F (Falso):
\begin{enumerate}[label=\Roman*., itemsep=\espacoAlt]
    \item Denominamos essa dinâmica de "rotação síncrona", o que significa que o tempo que a Lua leva para dar uma volta em torno de si mesma é exatamente igual ao tempo que ela leva para orbitar a Terra.
    \item A face oculta da Lua é conhecida como "lado escuro" porque nunca recebe a luz do Sol em nenhum momento do ciclo lunar de \textbf{29,5 dias}.
    \item Se a Lua parasse subitamente de girar em torno do seu próprio eixo, mas mantivesse sua órbita ao redor da Terra, os observadores terrestres passariam a enxergar toda a sua superfície ao longo de um mês.
\end{enumerate}
\par\vspace{\espacoQuestao}

\noindent\textcolor{CorLinha}{\rule{\linewidth}{0.5pt}} 
\par\vspace{\espacoPreQuestao}
\subsubsection*{15.} Os eclipses ocorrem quando há o alinhamento geométrico entre o Sol, a Terra e a Lua. O diagrama óptico detalha a formação dos cones de sombra durante um evento específico.

\begin{center}
% ==========================================
% CONTROLE INDIVIDUAL DESTA IMAGEM:
% ==========================================
\begin{adjustbox}{trim=0cm 0cm 0cm 0cm, clip, max width=\larguraTikz}
\begin{tikzpicture}
\definecolor{fundo}{HTML}{ECF0F1}
\definecolor{sol}{HTML}{F39C12}
\definecolor{terra}{HTML}{2980B9}
\definecolor{lua}{HTML}{7F8C8D}
\definecolor{umbra}{HTML}{2C3E50}
\definecolor{penumbra}{HTML}{95A5A6}

\fill[fundo, rounded corners=3pt] (-1, -3) rectangle (8, 3);

% LAYER 1: Cones de Sombra (Penumbra e Umbra)
\fill[penumbra, opacity=0.3] (0,2) -- (7,-1.5) -- (7,1.5) -- (0,-2) -- cycle;
\fill[umbra, opacity=0.6] (0,2) -- (4,0) -- (0,-2) -- cycle;

% LAYER 2: Astros
\fill[sol, drop shadow={opacity=0.3}] (-3,0) circle (2.5cm);
\node[font=\bfseries, text=white] at (-1.5, 0) {Sol};

\fill[lua, drop shadow={opacity=0.4}] (0,0) circle (0.4cm);
\node[font=\bfseries\tiny, fill=white, inner sep=1pt] at (0, 0.7) {Lua};

% LAYER 3: Regiões de Observação na Terra
\draw[thick] (4,0) arc (180:130:3cm);
\draw[thick] (4,0) arc (180:230:3cm);
\fill[terra, drop shadow={opacity=0.2}] (4.5,0) circle (1.2cm);

% Pontos de interesse
\fill[white] (3.4, 0) circle (1.5pt) node[right, font=\bfseries\footnotesize, text=black, fill=white, inner sep=1pt] {A};
\fill[white] (3.6, 0.8) circle (1.5pt) node[right, font=\bfseries\footnotesize, text=black, fill=white, inner sep=1pt] {B};
\fill[white] (4.5, -1) circle (1.5pt) node[right, font=\bfseries\footnotesize, text=black, fill=white, inner sep=1pt] {C};

\end{tikzpicture}
\end{adjustbox}
\end{center}

Com base no traçado óptico dos raios solares, determine e explique detalhadamente a diferença do fenômeno astronômico visualizado por um observador no ponto \textbf{A} e outro no ponto \textbf{B}, justificando através dos conceitos físicos de umbra e penumbra.
\par\vspace{\espacoQuestao}

\noindent\textcolor{CorLinha}{\rule{\linewidth}{0.5pt}} 
\par\vspace{\espacoPreQuestao}
\subsubsection*{16.} Durante uma prova de recuperação, a estudante Júlia tentou justificar a raridade dos eclipses escrevendo o seguinte argumento na prova:

\begin{boxresponda}
\textbf{Responda:} "Como a Lua cruza a frente da Terra a cada 29,5 dias na Fase Nova, nós deveríamos ter um eclipse solar todos os meses. O motivo disso não acontecer é que a velocidade da Lua é inconstante. Às vezes ela passa tão rápido na frente do Sol que a sombra não tem tempo de atingir a Terra."\\
\\
O argumento de Júlia contém uma falha conceitual grave envolvendo a mecânica orbital. Diagnostique esse erro e apresente a verdadeira justificativa astronômica (relacionada ao plano da órbita) que explica por que não temos eclipses mensais.
\end{boxresponda}
\par\vspace{\espacoQuestao}

\noindent\textcolor{CorLinha}{\rule{\linewidth}{0.5pt}} 
\par\vspace{\espacoPreQuestao}
\subsubsection*{17.} (ENEM - Adaptada) A força gravitacional que a Lua e o Sol exercem sobre as massas de água do nosso planeta é responsável pelo fenômeno das marés. Sabe-se que as "marés de sizígia" (marés vivas, de maior amplitude) ocorrem quando Sol, Terra e Lua estão perfeitamente alinhados, somando suas forças gravitacionais. Esse alinhamento máximo coincide estritamente com quais fases lunares?
\begin{enumerate}[label=\alph*), itemsep=\espacoAlt]
    \item Fases Crescente e Minguante.
    \item Fases Nova e Cheia.
    \item Apenas na fase Cheia, devido ao aumento da reflexão de luz solar.
    \item Apenas na fase Nova, quando a Lua cruza o Trópico de Capricórnio.
    \item Nas fases de Quarto Crescente, em conjunção direta com os equinócios.
\end{enumerate}
\begin{boxdica}
\textbf{Refutação Obrigatória:} Escolha a alternativa correta e escreva um parágrafo provando fisicamente por que a alternativa "a)" está profundamente incorreta no que diz respeito aos vetores gravitacionais.
\end{boxdica}
\par\vspace{\espacoQuestao}

\noindent\textcolor{CorLinha}{\rule{\linewidth}{0.5pt}} 
\par\vspace{\espacoPreQuestao}
\subsubsection*{18.} A órbita da Lua ao redor da Terra não é um círculo perfeito, mas sim uma elipse. Isso significa que a distância entre os dois astros varia, possuindo um ponto de maior proximidade (perigeu) e um de maior afastamento (apogeu). Durante um alinhamento para um eclipse solar, se a Lua estiver no seu apogeu, ela não conseguirá encobrir todo o disco solar. Como é classificado e visualmente caracterizado o eclipse solar que ocorre nessas condições específicas de distanciamento?
\par\vspace{\espacoQuestao}

\noindent\textcolor{CorLinha}{\rule{\linewidth}{0.5pt}} 
\par\vspace{\espacoPreQuestao}
\subsubsection*{19.} O eclipse lunar total ocorre quando a Lua mergulha completamente no cone de umbra projetado pela Terra. Observe a sequência da trajetória de inserção lunar simulada em um plano de perfil.

\begin{center}
% ==========================================
% CONTROLE INDIVIDUAL DESTA IMAGEM:
% ==========================================
\begin{adjustbox}{trim=0cm 0cm 0cm 0cm, clip, max width=\larguraTikz}
\begin{tikzpicture}
\definecolor{ceu}{HTML}{1A1A24}
\definecolor{umbra}{HTML}{34495E}
\definecolor{penumbra}{HTML}{5D6D7E}
\definecolor{luabranca}{HTML}{ECF0F1}
\definecolor{luasangue}{HTML}{C0392B}

\fill[ceu, rounded corners=2pt, drop shadow={opacity=0.3}] (0, -2) rectangle (8, 4);

% Cones projetados transversalmente
\fill[penumbra, opacity=0.5] (4,1) circle (2.5cm);
\fill[umbra, opacity=0.8] (4,1) circle (1.2cm);
\node[fill=white, text=black, font=\bfseries\footnotesize, inner sep=1pt] at (4, 1) {Umbra da Terra};
\node[fill=white, text=black, font=\bfseries\footnotesize, inner sep=1pt] at (4, 3) {Penumbra};

% Trajetória
\draw[dashed, thick, white] (0.5, -0.5) -- (7.5, 2.5);

% Luas
\fill[luabranca] (1.2, -0.2) circle (0.3cm) node[above=4pt, text=white, font=\bfseries\tiny] {T1};
\fill[luasangue] (4, 1) circle (0.3cm) node[below=4pt, text=white, font=\bfseries\tiny] {T2};
\fill[luabranca] (6.8, 2.2) circle (0.3cm) node[above=4pt, text=white, font=\bfseries\tiny] {T3};

\end{tikzpicture}
\end{adjustbox}
\end{center}

Durante o instante \textbf{T2}, quando a Lua está totalmente imersa na umbra da Terra (área teórica de escuridão total), ela não desaparece completamente do céu noturno, mas adquire uma coloração avermelhada característica ("Lua de Sangue"). Explique analiticamente o fenômeno de refração atmosférica terrestre que permite que a luz atinja a Lua nesse instante.
\par\vspace{\espacoQuestao}

\noindent\textcolor{CorLinha}{\rule{\linewidth}{0.5pt}} 
\par\vspace{\espacoPreQuestao}
\subsubsection*{20.} A observação astronômica depende profundamente do referencial adotado. Imagine que uma base de exploração foi instalada na superfície da Lua, na face voltada para nós, e um astronauta observa o espaço sideral.

\begin{center}
% ==========================================
% CONTROLE INDIVIDUAL DESTA IMAGEM:
% ==========================================
\begin{adjustbox}{trim=0cm 0cm 0cm 0cm, clip, max width=\larguraTikz}
\begin{tikzpicture}
\definecolor{fundo}{HTML}{212F3D}
\definecolor{terra}{HTML}{3498DB}
\definecolor{lua}{HTML}{7F8C8D}
\definecolor{sol}{HTML}{F1C40F}

\fill[fundo, rounded corners=4pt] (-2,-2) rectangle (7,2);

% Sol
\fill[sol, drop shadow={opacity=0.3}] (-3,0) circle (1.5cm);
\draw[->, thick, sol] (-1.2, 0.5) -- (0, 0.5);
\draw[->, thick, sol] (-1.2, -0.5) -- (0, -0.5);
\node[fill=fundo, text=sol, font=\bfseries] at (-1.5, 1.2) {Sol};

% Terra
\fill[terra] (1.5,0) circle (0.8cm);
\node[fill=white, text=terra, font=\bfseries\tiny, inner sep=1pt] at (1.5, 1) {Terra};
% Sombreamento Terra
\begin{scope}
    \clip (1.5,-0.8) rectangle (2.3,0.8);
    \fill[black, opacity=0.7] (1.5,0) circle (0.8cm);
\end{scope}

% Lua e Base
\fill[lua] (5,0) circle (0.4cm);
\node[fill=white, text=black, font=\bfseries\tiny, inner sep=1pt] at (5, 0.6) {Lua};
% Sombreamento Lua
\begin{scope}
    \clip (5,-0.4) rectangle (5.4,0.4);
    \fill[black, opacity=0.7] (5,0) circle (0.4cm);
\end{scope}

% Observador
\draw[->, thick, white, dashed] (4.6, 0) -- (2.5, 0);
\node[fill=fundo, text=white, font=\bfseries\tiny, align=center] at (3.5, -0.5) {Visão do\\Astronauta};

\end{tikzpicture}
\end{adjustbox}
\end{center}

Baseando-se no posicionamento esquematizado (em que a Terra se encontra perfeitamente entre o Sol e a Lua), julgue os cenários abaixo e explique o raciocínio:
Qual seria a fase da Lua vista pelos humanos na Terra nesse exato momento e, simultaneamente, qual seria a aparência da "fase da Terra" visualizada pelo astronauta na base lunar?
\par\vspace{\espacoQuestao}

\noindent\textcolor{CorLinha}{\rule{\linewidth}{0.5pt}} 
\par\vspace{\espacoPreQuestao}
\subsubsection*{21.} (Estilo VUNESP) Em documentários científicos, é comum escutar que a gravidade na superfície lunar é cerca de um sexto ($ \mfrac{1}{6} $) da gravidade na superfície terrestre.
\begin{enumerate}[label=\alph*), itemsep=\espacoAlt]
    \item A massa da Lua é extremamente menor, o que a impede de reter moléculas gasosas, resultando na ausência de atmosfera.
    \item O raio da Lua é um sexto do raio da Terra, resultando em uma proporção geométrica perfeita de atração.
    \item A rotação lenta da Lua neutraliza a força centrífuga, diminuindo o peso aparente dos corpos em sua superfície.
    \item A Lua sofre influência magnética direta do Sol, anulando parte de sua própria atração gravitacional.
    \item A composição do solo lunar não possui metais pesados como o ferro, impossibilitando a geração de gravidade intensa.
\end{enumerate}
\begin{boxdica}
\textbf{Refutação Obrigatória:} Escolha a alternativa correta que explica o fenômeno atrelado à baixa gravidade lunar e escreva um diagnóstico justificando por que a alternativa "c)" trata-se de um conceito físico incorreto.
\end{boxdica}
\par\vspace{\espacoQuestao}

\noindent\textcolor{CorLinha}{\rule{\linewidth}{0.5pt}} 
\par\vspace{\espacoPreQuestao}
\subsubsection*{22.} Um fórum de internet sobre astronomia debatia a iluminação do nosso satélite natural, e um dos usuários publicou o seguinte comentário:

\begin{boxresponda}
\textbf{Responda:} "As fases da Lua (como a Crescente e a Minguante) ocorrem por causa da sombra da Terra. Conforme a Lua gira, a Terra entra na frente do Sol e projeta uma sombra parcial nela. Quando a sombra da Terra cobre a Lua inteira, nós chamamos de Lua Nova."\\
\\
O comentário confunde dois fenômenos astronômicos completamente distintos. Diagnostique a confusão lógica do usuário, separando o que causa as Fases da Lua do que realmente representa a projeção da sombra da Terra.
\end{boxresponda}
\par\vspace{\espacoQuestao}

\noindent\textcolor{CorLinha}{\rule{\linewidth}{0.5pt}} 
\par\vspace{\espacoPreQuestao}
\subsubsection*{23.} O mês sinódico (tempo para a Lua retornar à mesma fase, como de Nova para Nova) dura cerca de \textbf{29,5 dias}. No entanto, o mês sideral (tempo que a Lua leva para completar uma volta de $360^\circ$ ao redor da Terra em relação às estrelas de fundo) dura apenas cerca de \textbf{27,3 dias}. Utilizando a geometria da translação da Terra ao redor do Sol, explique analiticamente por que o ciclo de fases lunares exige esses \textbf{2,2 dias} adicionais para ser completado.
\par\vspace{\espacoQuestao}

\noindent\textcolor{CorLinha}{\rule{\linewidth}{0.5pt}} 
\par\vspace{\espacoPreQuestao}
\subsubsection*{24.} O plano da órbita lunar apresenta uma inclinação de aproximadamente $5^\circ$ em relação ao plano da órbita terrestre (a Eclíptica). Desenhe mentalmente ou crie um esquema escrito provando que, sem essa inclinação angular de $5^\circ$, os habitantes da Terra vivenciariam um evento astronômico drástico e repetitivo no céu a cada fase de Lua Cheia. Que evento seria esse?
\par\vspace{\espacoQuestao}

\noindent\textcolor{CorLinha}{\rule{\linewidth}{0.5pt}} 
\par\vspace{\espacoPreQuestao}
\subsubsection*{25.} O Programa Artemis, da NASA, planeja estabelecer presença sustentável no Polo Sul lunar. Um dos maiores desafios de engenharia para a sobrevivência das missões é o controle térmico e de energia solar, visto que a duração do "dia" e da "noite" na Lua é ditada exclusivamente pela sua rotação síncrona. Sabendo que o período de rotação da Lua tem a duração de aproximadamente um mês terrestre, demonstre matematicamente a duração aproximada, em dias terrestres, de um período contínuo de escuridão total (a "noite lunar") enfrentado por uma base posicionada na região equatorial da Lua.
\par\vspace{\espacoQuestao}

\end{multicols*}
\end{document}